% sec-05: authority inside the trial.  Table 15.
\section{Authority Inside the Robotic Phase 1}

The same principle that places the company below every structure it writes to
also governs the trial itself. The chief executive holds all of the sponsor, so
every trigger in the federal financial disclosure rule would fire if the same
person held an investigator's role as well. The site package therefore places every decision an
interested party could move with a role at the site, and leaves the company only
the signature a sponsor cannot delegate \cite{movein2026}. Table 15 lists the
decisions and who holds each.

\subsection{The trial those decisions govern}

The deposited protocol is an open-label, single-arm, 3+3 dose escalation of
perioperative daraxonrasib at 160, 220 and 300 mg, with a 28-day dose-limiting
toxicity window, in up to 18 treated participants with resectable or
borderline-resectable KRAS-mutated pancreatic ductal adenocarcinoma
\cite{onpremwhippl}. The procedure is a staged eight-arm robotic
pancreaticoduodenectomy with 56 degrees of freedom and 640 sensor channels, force
limits of 3 N per arm and 18 N cumulative, and a five-vessel no-fly gate. An
on-premises language model is confined to advisory output, with no write
credential to data capture and no route to the robot control network
\cite{phase1ind}. The investigational new drug application it would be filed
under follows 21 CFR Part 312 \cite{ecfr2026part312}.

\begin{apptable}
\begin{tabularx}{\textwidth}{>{\raggedright\arraybackslash}p{5.3cm} >{\raggedright\arraybackslash}p{4.5cm} >{\raggedright\arraybackslash}X}
\headrow \thead{Decision} & \thead{Held by} & \thead{Why it sits there}\\
\midrule
Dose assignment in the 3+3 rule & Biostatistician & The only discretionary step \\
Dose-limiting toxicity calls & Site and sub-investigator & Decides if escalation continues \\
Analysis plan, database lock & Biostatistician & Decides what the trial reports \\
Consent and eligibility & Investigator, coordinator & Participant safety starts here \\
Use of the model's advice & The operating surgeon & The model is advisory only \\
Emergency stop chain & Site safety officer & 3 ms arm stop; 500 ms system stop \\
Regulatory signature & Chief executive & A sponsor duty, never delegated \\
\end{tabularx}
\end{apptable}
\vspace{-0.60cm}
\tabcap{Table 15. The seven decisions inside the robotic Phase 1 and the role that holds\\
each; the company keeps only the regulatory signature a sponsor cannot delegate.}

\subsection{The evidence behind the choice of agent}

The agent was chosen on a dated, quantitative record. In August 2025 the
company's ten-arm simulation returned a median overall survival of 12.8 months
against 5.4 for chemotherapy \cite{qspmetpancre}. In May 2026 the developer's
Phase 3 trial reported 13.2 months against 6.6 in the RAS G12 population of
previously treated metastatic disease \cite{rasolute302}. The two figures are
close, and the company describes that closeness as a chronology rather than a
validation, because the populations, regimens and sample sizes differ. Pancreatic
ductal adenocarcinoma remains the disease with the widest gap between incidence
and survival in the United States \cite{siegel2025statistics}, and the standard
of care after resection remains combination chemotherapy
\cite{conroy2018folfirinox}.

\subsection{The money the decisions would spend}

The program's direct cost is \$700,000 a year for five years, of which \$521,000
is personnel across 3.95 full-time equivalents \cite{movein2026}. The SBIR route
of \$306,000 in Phase I and \$1,300,000 in Phase II totals \$1,606,000, and the
capitalization plan places \$5,900,000 of private capital behind a firewall
against the \$2,104,000 delta the route does not buy \cite{capitalplan2026}. The
company's comparable virtual trial work is \textbf{projected} at \$36,330 per
run, a figure that describes simulation work already deposited and not the cost
of the clinical trial.
